% Template for your submission.pdf. Enter your team name below, write your answers after "Answer:", and compile
% this file to create submission.pdf.
\documentclass[addpoints]{exam}

% ---------- Fill in ----------
\newcommand{\teamname}{my-team}  % your team name, as registered
% ---------- End of the part to fill in ----------

\newcommand{\assignmentnumber}{4}

% ---------- Layout of the assignment sheets (from the sheets' macros.tex) ----------
\usepackage{amsmath, amsfonts, amssymb}
\usepackage{graphicx}
\usepackage[super]{nth}
\usepackage{parskip}
\usepackage[hyperfootnotes=false,colorlinks]{hyperref}
\usepackage[utf8]{inputenc}
\usepackage[T1]{fontenc}
\usepackage{lmodern}
\usepackage[protrusion=false,expansion=false]{microtype}
\DisableLigatures{encoding = T1, family = tt*}
\usepackage{enumitem}
\pagestyle{headandfoot}

\newcommand{\lecture}{Deep Learning}
\newcommand{\semester}{WS 2026/2027}
\newcommand{\assignment}[1]{\nth{#1} Assignment}

\newcommand{\formatqwithpoints}{
\qformat{\large \thequestion. [\thepoints] \textbf{\thequestiontitle}\hfill\vrule depth \spacingbelowquestion width 0pt}
\bonusqformat{\large \thequestion. [\thepoints] \textbf{\thequestiontitle}\hfill\vrule depth \spacingbelowquestion width 0pt}
}
\newcommand{\formatqwithoutpoints}{
\qformat{\large \thequestion. \textbf{\thequestiontitle}\hfill\vrule depth \spacingbelowquestion width 0pt}
\bonusqformat{\large \thequestion. \textbf{\thequestiontitle}\hfill\vrule depth \spacingbelowquestion width 0pt}
}
\newcommand{\spacingbelowquestion}{0.85em}
\formatqwithpoints
\bonuspointpoints{bonus point}{bonus points}
\renewcommand{\thepartno}{\arabic{partno}}
\renewcommand\partlabel{\thepartno)}
\bracketedpoints

\setlength\parindent{0em}
\setlength{\parskip}{1.0em}
\renewcommand{\baselinestretch}{1.0}
\renewcommand{\questionshook}{\setlength{\itemsep}{0.7em}
    \setlength{\leftmargin}{0pt}\setlength{\labelwidth}{-\labelsep}}
\renewcommand{\partshook}{\setlength{\itemsep}{0em}
    \setlength{\leftmargin}{1em}
    \setlength{\labelwidth}{-\labelsep}
}
\setlist[itemize]{itemsep=0em,parsep=0em}
\newcommand{\myhrule}{\rule{\textwidth}{.5pt}}

\newcommand{\titledq}[2]{\titledquestion{#1}[#2]\setlength{\parskip}{1.0em}}
\newcommand{\titledqwithparts}[1]{
    \large \formatqwithoutpoints
    \titledquestion{#1}[]
    \formatqwithpoints
    \setlength{\parskip}{1.0em}
    \normalsize
}
\newcommand{\code}[1]{\texttt{#1}}
\newcommand{\fileref}[1]{\texttt{#1}}
\newcommand{\tbold}[1]{\textbf{#1}}
\newcommand{\hint}{\tbold{Hint:}\ }
\newcommand{\answer}{\tbold{Answer:}\ }

% questions and parts with the number, title and points they have on the sheet
% (the exam class keeps counting 1, 2, 3 internally; only the printed numbers are the sheet's)
\newcommand{\sheetnumber}{}
\newcommand{\sheetpartnumber}{}
\renewcommand{\thequestion}{\sheetnumber}
\renewcommand{\thepartno}{\sheetpartnumber}
\newcommand{\sheetq}[3]{\renewcommand{\sheetnumber}{#1}\titledq{#2}{#3}}
\newcommand{\sheetqwithparts}[2]{\renewcommand{\sheetnumber}{#1}\titledqwithparts{#2}}
\newcommand{\sheetpart}[3]{\renewcommand{\sheetpartnumber}{#1}\part[#2] \tbold{#3}}

\firstpageheader{Team: \teamname}{{\textbf{\lecture}}\\ Submission for the \assignment{\assignmentnumber}}{\semester}
\firstpagefooter{}{\thepage}{}
\runningheader{Team: \teamname}{Submission for the \assignment{\assignmentnumber}}{\semester}
\runningfooter{}{\thepage}{}
\headrule
% ---------- End of the layout ----------

\newcommand{\addtemplatenotes}{
\setlength{\parskip}{0.0em}
\small{
\tbold{About this template:} This template contains the questions from the assignment sheet that require a
written answer, with the original numbering and titles. While you may submit a handwritten scan instead, we
strongly encourage using this \LaTeX{} template. \LaTeX{} is widely used for scientific papers and theses,
especially for mathematical notation, so practicing it here is good preparation for your future research and
thesis.
\begin{itemize}
    \item \tbold{Using this template:} Write your answer after \tbold{Answer:} below each question and compile
    \code{submission.tex} to create \code{submission.pdf}.
    \item \tbold{Showing your work:} Always provide your intermediate steps and reasoning, not just the final
    result. For mathematical questions (derivations, proofs or calculations), show each step. For conceptual
    questions, briefly explain the reasoning behind your answer.
    \item \tbold{Team name \& anonymity:} Enter your team name at the top of \code{submission.tex}. Do not add
    any other personal information, such as GitHub usernames, real names, email addresses or matriculation
    numbers.
    \item \tbold{How to hand in:} Commit and push \code{submission.pdf} to your team's repository together with
    \code{submission.tex} and all files that it uses (for example, images and \code{.bib} files).
\end{itemize}
\myhrule%
}%
\setlength{\parskip}{1.0em}%
}

\begin{document}
    \addtemplatenotes

    \begin{questions}

        \sheetqwithparts{1}{Pen and Paper tasks}

        \begin{parts}
            \sheetpart{1}{2}{}Perform two update steps of SGD with momentum to the weights
            $\mathbf{w} = \begin{bmatrix} w_1 \\ w_2 \end{bmatrix}$ of the linear function $\mathbf{\hat y} = X\mathbf{w}$
            optimizing the mean squared error (MSE). As in exercises 2 and 3, $X$ contains one data point per row. The
            loss for one sample is given by $(\hat y - y)^2$, and, as for minibatch SGD in the lecture (Stochastic
            Gradient Descent (SGD) \& Minibatch Sizes), the per-sample losses are summed and divided by the batch size.
            Use the momentum update from the lecture: $\mathbf{v} \leftarrow \beta \mathbf{v} - \alpha \mathbf{\hat g}$,
            $\mathbf{w} \leftarrow \mathbf{w} + \mathbf{v}$, where $\mathbf{\hat g}$ is the gradient of the loss with
            respect to $\mathbf{w}$.

            Given: the initial state $\mathbf{w} = \begin{bmatrix} 1 \\ 0 \end{bmatrix}$,
            velocity $\mathbf{v} = \begin{bmatrix} 0 \\ 0 \end{bmatrix}$,
            data $X = \begin{bmatrix} 2 & -1 \\ -1 & 3 \end{bmatrix}$,
            target $\mathbf{y} = \begin{bmatrix} 3 \\ 1 \end{bmatrix}$,
            learning rate $\alpha = 0.2$, momentum strength $\beta = 0.8$ and batch size $B = 2$.

            \answer Your answer.
            % For math over several lines, uncomment and fill in:
            % \begin{align*}
            %     \mathbf{\hat y} &= X\mathbf{w} \\
            %     L &= \frac{1}{B} \sum_{i=1}^{B} (\hat y_i - y_i)^2
            % \end{align*}

            \sheetpart{2}{2}{}Prove that in Adam the bias-corrected first moment estimate $\hat{s}$ (with $s$
            initialized to $0$) is equal to the gradient in every step when that gradient remains constant.

            \hint This can be done using mathematical induction.

            \answer Your answer.
        \end{parts}

        \sheetqwithparts{2}{Coding Tasks}

        \begin{parts}
            \sheetpart{3}{2}{}Compare Adam vs SGD vs SGD with Momentum.

            Create two plots, one with the training loss curves and one with the training accuracy curves of the
            three models (function \code{plot\_learning\_curves} in file \fileref{lib/plot.py}).

            Run \fileref{plot\_trained\_models.py} and include the two plots here.

            \answer

            % Save the plots as image files, add them to your repository and uncomment, for example:
            % \includegraphics[width=0.8\textwidth]{learning_curves.png}

            \sheetpart{5}{1}{}Plot learning rate schedules.

            To verify your implementation, plot the learning rate schedules with number of epochs on the x-axis and
            learning rate on the y-axis (function \code{plot\_lr\_schedules} in file \fileref{lib/plot.py}).

            Run \fileref{plot\_lr\_schedules.py} and include the plot here.

            \answer

            % \includegraphics[width=0.6\textwidth]{lr_schedules.png}
        \end{parts}

        \sheetqwithparts{3}{Experiments}

        \begin{parts}
            \sheetpart{1}{1}{}Run \fileref{plot\_ill\_conditioned.py} for the following cases and include all the
            generated plots:
            \begin{itemize}
                \item Condition number = 1, momentum = 0.0, preconditioning = False
                \item Condition number = 8, momentum = 0.0, preconditioning = False
                \item Condition number = 8, momentum = 0.35, preconditioning = False
                \item Condition number = 8, momentum = 0.0, preconditioning = True
            \end{itemize}

            \answer

            % \includegraphics[width=0.44\textwidth]{cond1_momentum0.png}
            % \includegraphics[width=0.44\textwidth]{cond8_momentum0.png}

            % \includegraphics[width=0.44\textwidth]{cond8_momentum035.png}
            % \includegraphics[width=0.44\textwidth]{cond8_momentum0_precond.png}

            Explain why the second experiment causes zigzag behaviour and why momentum and preconditioning can
            mitigate this issue.

            \answer Your answer.
        \end{parts}

    \end{questions}
\end{document}
